\section{The Toric Coherent--Constructible Correspondence}

We now state the toric coherent--constructible correspondence. Let \(X_\Sigma\)
be a smooth toric variety with acting torus \(T\), and let
\(\Lambda_\Sigma\subset T^*M_{\mathbb R}\) be the FLTZ skeleton associated to
the fan \(\Sigma\). The coherent side of the compact form of the correspondence
is the perfect equivariant category \(\operatorname{Perf}_T(X_\Sigma)\). In the
proof, we first work with the theta-generated equivariant quasi-coherent
category and then pass to compact objects, using
\cref{thm:compact-form-of-the-coherent-side}. The constructible side is the dg
category of constructible sheaves on \(M_{\mathbb R}\) with singular support
contained in \(\Lambda_\Sigma\).



\begin{theorem}[Toric coherent--constructible correspondence]\label{thm:toric-ccc}
Let \(X_\Sigma\) be a smooth toric variety. There is a quasi-equivalence of
dg categories
\[
\kappa_\Sigma:
\operatorname{Perf}_T(X_\Sigma)
\xrightarrow{\sim}
Sh_{\Lambda_\Sigma}(M_{\mathbb R}).
\]
Equivalently, on homotopy categories one obtains an equivalence
\[
D^b(\operatorname{Coh}_T(X_\Sigma))
\simeq
H^0\bigl(Sh_{\Lambda_\Sigma}(M_{\mathbb R})\bigr).
\]
\end{theorem}

The proof will be given by comparing the theta generators on the two sides.
The objects \(\Theta'(\sigma,\chi)\) and \(\Theta(\sigma,\chi)\) are indexed
by the same data \((\sigma,\chi)\), and the functor \(\kappa_\Sigma\) is
characterized by
\[
\kappa_\Sigma\bigl(\Theta'(\sigma,\chi)\bigr)=\Theta(\sigma,\chi).
\]
It remains to compare the Hom-complexes between these generators and then
apply the generator criterion for quasi-equivalences.


\subsection{The Correspondence of Generators}

The theta objects on the two sides are indexed by the same combinatorial
data. More precisely, for each cone \(\sigma\in\Sigma\), let
\(q_\sigma:M\to M_\sigma\) be the quotient map.

\begin{lemma}
Let \(\sigma\in\Sigma\), and let \(\chi,\chi'\in M\). If
\(q_\sigma(\chi)=q_\sigma(\chi')\), then
\(\Theta'(\sigma,\chi)\simeq \Theta'(\sigma,\chi')\).
Hence the coherent-side theta object depends only on the class
\(\bar\chi\in M_\sigma\).
\end{lemma}

\begin{proof}
The condition \(q_\sigma(\chi)=q_\sigma(\chi')\) means that
\(\chi'-\chi\in\sigma^\perp\cap M\). Put \(m_0=\chi'-\chi\). Then the
monomial \(\chi^{m_0}\) is an invertible regular function on \(U_\sigma\).
Multiplication by \(\chi^{m_0}\) gives a \(T\)-equivariant isomorphism
\(\mathcal O_{U_\sigma}(\chi)\simeq\mathcal O_{U_\sigma}(\chi')\), and
pushing forward along \(i_\sigma:U_\sigma\hookrightarrow X_\Sigma\) gives
\(\Theta'(\sigma,\chi)\simeq\Theta'(\sigma,\chi')\).
\end{proof}

Thus we may write the coherent-side theta object as
\(\Theta'(\sigma,\bar\chi)\), where \(\bar\chi\in M_\sigma\). The
constructible theta object \(\Theta(\sigma,\bar\chi)\) is already indexed by
the same data. The correspondence of generators is therefore
\[
\Theta'(\sigma,\bar\chi)\longmapsto \Theta(\sigma,\bar\chi).
\]

The functor \(\kappa_\Sigma\), once constructed, will be characterized on
theta generators by
\[
\kappa_\Sigma\bigl(\Theta'(\sigma,\bar\chi)\bigr)
=
\Theta(\sigma,\bar\chi).
\]
The next step is to compare the Hom-complexes between these generators. This
comparison is the essential computation from which the CCC functor and the
quasi-equivalence will follow.
\subsection{Hom-Complex Computations}
We now compare the Hom-complexes between theta objects. For
\(\bar\chi\in M_\sigma\), choose a lift \(\chi\in M\) and write
\(P_{\sigma,\bar\chi}:=\chi+\sigma^\vee\). This subset is independent of the
chosen lift.


\begin{lemma}[Hom between polyhedral costandard sheaves]
\label{lem:hom-polyhedral-costandard-sheaves}
Let \(P,Q\subset M_{\mathbb R}\) be closed convex polyhedra, and let
\(i_P:P\hookrightarrow M_{\mathbb R}\) and
\(i_Q:Q\hookrightarrow M_{\mathbb R}\) be the inclusions. Then
\[
\mathrm{RHom}((i_P)_*\omega_P,(i_Q)_*\omega_Q)
\simeq
\begin{cases}
k, & P\subset Q,\\
0, & \text{otherwise}.
\end{cases}
\]
In the nonzero case, the complex is concentrated in degree zero.
\end{lemma}

\begin{proof}
By Verdier duality, \((i_P)_*\omega_P\simeq D(k_P)\). Hence
\[
\mathrm{RHom}((i_P)_*\omega_P,(i_Q)_*\omega_Q)
\simeq
\mathrm{RHom}(D(k_P),D(k_Q))
\simeq
\mathrm{RHom}(k_Q,k_P).
\]
For closed convex polyhedra, \(\mathrm{RHom}(k_Q,k_P)\) is computed by the
cohomology of the region \(P\cap Q\) with the restriction condition from \(Q\)
to \(P\). If \(P\subset Q\), this region is the convex polyhedron \(P\), hence
contractible, and the Hom-complex is \(k\) in degree zero. If
\(P\not\subset Q\), the corresponding restriction condition fails, and the
Hom-complex is acyclic. This gives the stated formula.
\end{proof}


\begin{proposition}[Hom-complex comparison]
\label{prop:hom-complex-comparison}
Let \(\sigma,\tau\in\Sigma\), \(\bar\chi\in M_\sigma\), and
\(\bar\psi\in M_\tau\). Then there is a natural quasi-isomorphism
\[
\mathrm{RHom}_T\bigl(\Theta'(\sigma,\bar\chi),
\Theta'(\tau,\bar\psi)\bigr)
\simeq
\mathrm{RHom}\bigl(\Theta(\sigma,\bar\chi),
\Theta(\tau,\bar\psi)\bigr).
\]
More explicitly, after choosing lifts \(\chi,\psi\in M\), both sides are
quasi-isomorphic to
\[
\begin{cases}
k, & P_{\sigma,\bar\chi}\subset P_{\tau,\bar\psi},\\
0, & \text{otherwise}.
\end{cases}
\]
In the nonzero case, the complex is concentrated in degree zero.
\end{proposition}

\begin{proof}
We compute the two sides separately.

On the coherent side, write
\[
\Theta'(\sigma,\bar\chi)=(i_\sigma)_*\mathcal O_{U_\sigma}(\chi),
\qquad
\Theta'(\tau,\bar\psi)=(i_\tau)_*\mathcal O_{U_\tau}(\psi).
\]
Since \(U_\sigma\cap U_\tau=U_{\sigma\cap\tau}\) is again an affine toric
chart, the Hom-complex is computed on affine charts by the corresponding
graded module calculation. A \(T\)-equivariant morphism between the two
equivariant line bundles is given by multiplication by a homogeneous monomial.
With our linearization convention, the possible monomial has weight
\(\chi-\psi\).

This monomial defines a regular map on the relevant affine toric chart exactly
when the corresponding character lies in the appropriate semigroup. In
polyhedral terms, this condition is equivalent to
\[
P_{\sigma,\bar\chi}\subset P_{\tau,\bar\psi}.
\]
If this inclusion holds, the space of equivariant morphisms is one-dimensional,
generated by this monomial. If the inclusion fails, there is no equivariant
regular monomial map. Since the relevant toric intersections are affine, there
are no higher derived Hom groups. Hence
\[
\mathrm{RHom}_T\bigl(\Theta'(\sigma,\bar\chi),
\Theta'(\tau,\bar\psi)\bigr)
\simeq
\begin{cases}
k, & P_{\sigma,\bar\chi}\subset P_{\tau,\bar\psi},\\
0, & \text{otherwise}.
\end{cases}
\]

On the constructible side, the theta objects are costandard sheaves on the
closed convex polyhedra \(P_{\sigma,\bar\chi}\) and \(P_{\tau,\bar\psi}\):
\[
\Theta(\sigma,\bar\chi)
=
(i_{\sigma,\bar\chi})_*\omega_{P_{\sigma,\bar\chi}},
\qquad
\Theta(\tau,\bar\psi)
=
(i_{\tau,\bar\psi})_*\omega_{P_{\tau,\bar\psi}}.
\]
By \cref{lem:hom-polyhedral-costandard-sheaves}, we get
\[
\mathrm{RHom}\bigl(\Theta(\sigma,\bar\chi),
\Theta(\tau,\bar\psi)\bigr)
\simeq
\begin{cases}
k, & P_{\sigma,\bar\chi}\subset P_{\tau,\bar\psi},\\
0, & \text{otherwise}.
\end{cases}
\]

Therefore the coherent and constructible Hom-complexes are naturally
quasi-isomorphic. The comparison is compatible with composition: if
\[
P_{\sigma,\bar\chi}\subset P_{\tau,\bar\psi}
\subset P_{\nu,\bar\eta},
\]
then the generator of the first Hom-complex composes with the generator of the
second to give the generator of the third. If one of the inclusions fails, the
corresponding Hom-complex is zero. Hence the quasi-isomorphisms above are
compatible with composition.
\end{proof}



\subsection{The CCC Functor}
\label{subsec:ccc-functor}

We now construct the theta-generated form of the CCC functor. Let
\(\mathcal A'_\Sigma\) be the full dg subcategory of
\(D(\operatorname{QCoh}_T(X_\Sigma))\) whose objects are the coherent-side theta
objects \(\Theta'(\sigma,\bar\chi)\). Similarly, let \(\mathcal A_\Sigma\) be
the full dg subcategory of \(\operatorname{Sh}_{\Lambda_\Sigma}(M_{\mathbb R})\)
whose objects are the constructible theta objects \(\Theta(\sigma,\bar\chi)\).

Define a map on objects by
\[
\Theta'(\sigma,\bar\chi)\longmapsto \Theta(\sigma,\bar\chi).
\]
By the Hom-complex comparison in \cref{prop:hom-complex-comparison}, for every
pair of indices \((\sigma,\bar\chi)\) and \((\tau,\bar\psi)\), there is a natural
quasi-isomorphism
\[
\operatorname{RHom}_T(\Theta'(\sigma,\bar\chi),\Theta'(\tau,\bar\psi))
\simeq
\operatorname{RHom}(\Theta(\sigma,\bar\chi),\Theta(\tau,\bar\psi)).
\]
These quasi-isomorphisms are compatible with composition. Hence the assignment
on theta objects extends to a dg functor
\[
\kappa^0_\Sigma:\mathcal A'_\Sigma\longrightarrow \mathcal A_\Sigma .
\]

To obtain the theta-generated form of the CCC functor, we extend
\(\kappa^0_\Sigma\) to the triangulated dg subcategories generated by the theta
objects. Since exact dg functors preserve shifts, cones, and finite direct sums,
the assignment extends to an exact dg functor
\[
\kappa_\Sigma:
\left\langle
\Theta'(\sigma,\bar\chi)\mid \sigma\in\Sigma,\ \bar\chi\in M_\sigma
\right\rangle
\longrightarrow
\left\langle
\Theta(\sigma,\bar\chi)\mid \sigma\in\Sigma,\ \bar\chi\in M_\sigma
\right\rangle .
\]
Equivalently, \(\kappa_\Sigma\) is the exact dg functor characterized by
\[
\kappa_\Sigma(\Theta'(\sigma,\bar\chi))=\Theta(\sigma,\bar\chi)
\]
on theta generators, together with the Hom-complex identifications from
\cref{prop:hom-complex-comparison}.

At this stage, \(\kappa_\Sigma\) is quasi-fully faithful on the theta
generators. The compact coherent side will be obtained by passing to compact
objects, as in \cref{thm:compact-form-of-the-coherent-side}. The remaining step
is to use the generation results to prove that the resulting functor is a
quasi-equivalence.

\subsection{Proof of the Quasi-Equivalence}
\label{subsec:proof-quasi-equivalence}

We now prove that the functor constructed above gives the toric CCC. By
construction, \(\kappa_\Sigma\) sends each coherent-side theta object to the
corresponding constructible theta sheaf:
\[
\kappa_\Sigma(\Theta'(\sigma,\bar\chi))=\Theta(\sigma,\bar\chi).
\]
Moreover, \cref{prop:hom-complex-comparison} shows that, for any
\(\sigma,\tau\in\Sigma\), \(\bar\chi\in M_\sigma\), and
\(\bar\psi\in M_\tau\), the induced map on Hom-complexes
\[
\operatorname{RHom}_T(\Theta'(\sigma,\bar\chi),\Theta'(\tau,\bar\psi))
\longrightarrow
\operatorname{RHom}(\Theta(\sigma,\bar\chi),\Theta(\tau,\bar\psi))
\]
is a quasi-isomorphism.

It remains to use the generation results. On the coherent side,
\cref{thm:coherent-theta-generation} says that the theta objects generate the
equivariant quasi-coherent derived category
\(D(\operatorname{QCoh}_T(X_\Sigma))\) as a localizing triangulated category.
On the constructible side, \cref{thm:fltz-generation} says that the theta
sheaves generate \(\operatorname{Sh}_{\Lambda_\Sigma}(M_{\mathbb R})\). Hence
the hypotheses of the generator criterion
\cref{thm:generator-criterion} are satisfied for the theta-generated
categories: the functor is exact, it identifies the Hom-complexes between the
generators, and the corresponding theta objects generate the two sides.

Therefore \(\kappa_\Sigma\) gives a quasi-equivalence between the
theta-generated coherent and constructible categories. Passing to compact
objects on the coherent side, \cref{thm:compact-form-of-the-coherent-side}
identifies the compact part of \(D(\operatorname{QCoh}_T(X_\Sigma))\) with
\(\operatorname{Perf}_T(X_\Sigma)\). Thus we obtain the perfect form of the
toric coherent--constructible correspondence:
\[
\kappa_\Sigma:
\operatorname{Perf}_T(X_\Sigma)
\xrightarrow{\sim}
\operatorname{Sh}_{\Lambda_\Sigma}(M_{\mathbb R}).
\]

Since \(X_\Sigma\) is smooth, \(\operatorname{Perf}_T(X_\Sigma)\) is a dg
enhancement of \(D^b(\operatorname{Coh}_T(X_\Sigma))\). Hence, on homotopy
categories, one obtains
\[
D^b(\operatorname{Coh}_T(X_\Sigma))
\simeq
H^0\bigl(\operatorname{Sh}_{\Lambda_\Sigma}(M_{\mathbb R})\bigr).
\]
This proves the toric coherent--constructible correspondence.
